%! AP Statistics — Unit 7, Lesson 7.3 — Warm-Up
\documentclass[10pt]{article}
\usepackage{apstats-article}
\usepackage{apstats-boxes}
\newcommand{\TallMath}[1]{$\displaystyle #1\rule[-1.4em]{0pt}{3.2em}$}

\begin{document}

\pageheader{Unit 7, Lesson 7.3}{Warm-Up}
\namedateperiod

\vspace{0.15in}

From Lesson 7.2, you constructed this interval:\\[4pt]
\centerline{\textbf{95\% CI for mean sodium intake: $(2615, 3065)$ mg}}

\begin{enumerate}

  \item Three students wrote interpretations of this interval. Circle the
        \textbf{correct} interpretation and identify the error in each incorrect one.

        \begin{enumerate}[label=\Alph*.]
          \item ``There is a 95\% probability that the true mean sodium intake is
                between 2615 and 3065 mg.''

                Error (if any): \writelines{1}

          \item ``We are 95\% confident that the interval (2615, 3065) mg captures
                the true mean daily sodium intake for all college students.''

                Error (if any): \writelines{1}

          \item ``95\% of all college students have a daily sodium intake between
                2615 and 3065 mg.''

                Error (if any): \writelines{1}
        \end{enumerate}

  \item From Unit 6, a 95\% CI for a proportion $p$ was $(0.546, 0.734)$. The researcher
        claims $p = 0.80$. Does the interval provide convincing evidence against this claim?
        Explain in one sentence.

        \writelines{2}

  \item Preview: If the sample size in the sodium study were doubled (from $n = 20$ to
        $n = 40$, with the same $\bar x$ and $s$), would the margin of error increase,
        decrease, or stay the same? By approximately what factor?

        \writelines{2}

\end{enumerate}

\end{document}
